\src{175}
\result{PROPOSITION}{4} Let $A$ be a commutative ring with identity, and let
$(x_1,\ldots,x_p)$ be a sequence of elements of $A$ such that, for
$1\leq i\leq p$, the image of $x_i$ in the quotient of $A$ by the ideal
generated by $(x_1,\ldots,x_{i-1})$ is not a zero divisor. Let $J$ be the ideal
generated by the $x_i$. Then $J/J^{2}$ is a free $(A/J)$-module with basis the
canonical images of the $x_i$, the complex $K_{\underline{x}}$ is a free
resolution of $A/J$, and for every $A$-module $M$, the homomorphisms (3.1) in
degree $p$ are bijective. The same is true of the analogous homomorphisms
defined in arbitrary degrees $i$, provided that $J M=0$.

(The essential point, from which all the others follow, is the acyclicity of
$K_{\underline{x}}$, a well-known fact under the stated hypotheses.)

\result{COROLLARY}{1} With $A$ and $J$ as above, suppose in addition that $A$
is a regular affine algebra of dimension $n$ over a perfect field $k$, and that
$A/J$ is a regular affine algebra. Denote by $\Omega^{1}(A)$ and
$\Omega^{1}(A/J)$ the modules of Kähler differentials. Then there is a canonical
isomorphism compatible with localization,
\begin{equation}\tag{3.5}
  \operatorname{Ext}_{A}^{p}\bigl(\Omega^{n-p}(A/J),\Omega^{n}(A)\bigr)=A/J.
\end{equation}
Indeed, $\Omega^{n-p}(A/J)$ is a projective $(A/J)$-module of rank $1$, and likewise
$\Omega^{n}(A)$ is a projective $A$-module of rank $1$.\footnote{The printed text says
“rank $n$.” Both modules of top-degree forms have rank $1$, as is also required
by the identification with $\bigwedge^{p}(J/J^{2})$ below.} Thus the left-hand
side equals
\[
  \operatorname{Ext}_{A}^{p}(A/J,A)
  \otimes\Omega^{n-p}(A/J)'\otimes\Omega^{n}(A)
\]
(where the prime denotes the dual $(A/J)$-module). The tensor product of the
last two factors identifies with $\bigwedge^{p}(J/J^{2})$, so the whole expression
identifies with
\[
  \operatorname{Ext}_{A}^{p}\!\left(A/J,\bigwedge^{p}(J/J^{2})\right),
\]
and therefore, by the proposition, with
\[
  \operatorname{Hom}_{A}\!\left(\bigwedge^{p}J/J^{2},
  \bigwedge^{p}J/J^{2}\right),
\]
that is, with $A/J$. In particular,
$\operatorname{Ext}_{A}^{p}\bigl(\Omega^{n-p}(A/J),\Omega^{n}(A)\bigr)$ contains
a distinguished element corresponding to the identity element of $A/J$, called
the \emph{fundamental class of the ideal $J$ in $A$}. (It can in fact be defined
under considerably broader hypotheses.) Corollary 1 can be put in a more
geometric and more global form:

\result{COROLLARY}{2} Let $X$ be a nonsingular variety over an algebraically
closed field $k$, let $Y$ be a closed nonsingular subvariety of $X$, let
$\Osh_X$ be the sheaf of local rings of $X$, and let $\Osh_Y$ be the sheaf of
local rings of $Y$, regarded as a quotient sheaf of $\Osh_X$. Let $n$ be the
dimension of $X$ and $n-p$ that of $Y$.

